\section{Method}
\label{sec:method}
In this section, we introduce \textbf{ICO-GS}, a sparse-view reconstruction approach 
that enforces intrinsic consistency by coupling geometric regularization with 
appearance optimization. 
The overall framework is illustrated in \cref{fig:framework}. 
We first analyze the geometry-appearance discrepancy phenomenon under sparse-view 
settings (\cref{subsec:motivation}), then elaborate our robust geometric 
regularization (\cref{subsec:geometry}) and geometry-guided appearance 
optimization (\cref{subsec:appearance}). 
The complete optimization pipeline is described in \cref{subsec:optimization}.



\subsection{Motivation}
\label{subsec:motivation}

\begin{figure}[t]
  \vspace{-0.3cm}
  \begin{center}
     \includegraphics[trim={0cm 0cm 0cm 0cm},clip,width=1.0\linewidth]{figures/sparse2dense.pdf}
  \end{center}
  \vspace{-0.8cm}
  % \caption{ \textbf{Motivation.} }
  \caption{\textbf{Geometry-appearance discrepancy under sparse-view settings.}
  From top to bottom: RGB on training views, depth on training views, and RGB on test 
  views, rendered by 3D Gaussian Splatting~\cite{kerbl20233d} with varying training 
  view densities. With decreasing views, training-view appearance (top) remains 
  well-fitted, but depth quality (middle) collapses with noise and floaters
  due to insufficient multi-view constraints. This geometry-appearance discrepancy 
  leads to severe artifacts in novel-view rendering (bottom).
}
  % \label{fig:motivation}
  \label{fig:motivation}
  \vspace{-0.6cm}
  \end{figure}

  Despite success in dense-view reconstruction, 3DGS~\cite{kerbl20233d} severely 
  overfits under sparse observations. 
  We reveal a critical \textbf{geometry-appearance discrepancy}: as view count 
  decreases, appearance quality remains superficially high on training views while 
  geometric accuracy collapses, causing severe artifacts on novel views 
  (\cref{fig:motivation}).
  
  We identify two underlying deficiencies:
  
  \noindent\textbf{Insufficient Geometric Constraints.} 
  In dense-view settings, each 3D point is observed by multiple cameras (typically 10+), 
  providing strong multi-view constraints on Gaussian positions. 
  However, with sparse views, each point is visible in only 2--3 views, leading to 
  severe geometric ambiguity. 
  Mathematically, the photometric loss in \cite{kerbl20233d} only constrains the 
  \emph{projected} appearance of Gaussians, not their depth. 
  Consequently, Gaussians can be placed at \emph{any} position along the camera ray 
  while maintaining zero photometric error. 
  This depth ambiguity explains the noisy geometry in \cref{fig:motivation}: without 
  sufficient multi-view overlap, optimization lacks constraints to determine correct 
  3D positions.
    
  \noindent\textbf{Unreliable Appearance-Geometry Coupling.} 
  Moreover, the 3DGS representation does not inherently ensure geometry-appearance 
  consistency. 
  In principle, when a Gaussian is misplaced, the photometric loss should drive the 
  optimizer to correct its position. 
  However, we observe a problematic shortcut: instead of moving the Gaussian, the 
  optimizer adjusts its color and opacity to compensate for geometric errors, which 
  we term \emph{appearance compensation}.
  As shown in \cref{fig:motivation}, training views achieve high PSNR despite 
  severely incorrect geometry (noisy depth), demonstrating that appearance parameters 
  mask geometric mistakes. 
  This becomes particularly severe under sparse views, where weak geometric constraints 
  enable the model to overfit through appearance manipulation rather than learning 
  correct 3D structure.
  
  To address these deficiencies, we propose \textbf{ICO-GS}, which restores intrinsic 
  consistency by coupling robust geometric regularization (\cref{subsec:geometry}) 
  with geometry-guided appearance optimization (\cref{subsec:appearance}). 
  These components work synergistically to promote both geometry and appearance.
  


\subsection{Robust Geometric Regularization}
\label{subsec:geometry}
% ===== 开头：呼应前面的分析 =====
As identified in \cref{subsec:motivation}, sparse-view 3DGS suffers from insufficient geometric constraints. 
BinocularGS~\cite{han2024binocular} attempts to address this by enforcing stereo consistency via depth-warped virtual views. 
However, this approach suffers from a fundamental circular dependency: unreliable depth produces misaligned virtual views, which in turn provide corrupted supervision that further degrades geometry.

We break this loop via \emph{robust geometry regularization}, which warps pixels between training views according to rendered depth and penalizes inconsistencies. 
To handle illumination variations and occlusions that undermine naïve photometric matching, we introduce robust multi-view photometric consistency (\cref{subsubsec:mpc}) enhanced with edge-aware depth smoothness (\cref{subsubsec:smooth}).






\subsubsection{Multi-view Photometric Consistency}
\label{subsubsec:mpc}

% \noindent\textbf{Multi-view Photometric Consistency.}
Given $n$ sparse training views $\{I_i\}_{i=0}^{n-1}$, our goal is to regularize the rendered Gaussian depths $\{D_i\}_{i=0}^{n-1}$ via multi-view photometric consistency.
Take one reference view $I_{\text{0}}$ and its corresponding source views $\{I_j\}_{j=1}^{n-1}$ for example, we first render the depth map $D_{\text{0}}$ via alpha-blending rendering. 
Then for each pixel $p$ in the reference image $I_{\text{0}}$, its corresponding pixel $p'_{j}$ in the source images $\{I_j\}_{j=1}^{n-1}$ can be computed via:
\begin{equation}
  \vspace{-0.2cm}
  \label{eq:inverse_warp} 
    p'_{j} = KT_{0 \rightarrow j}(D_{0}(p) \cdot K^{-1} p),    
  % \vspace{-0.3cm}
\end{equation}
where $K, T$ denote the associated intrinsic and the relative transformation.
We enforce inverse warp from the source views to reference views to acquire the reconstructed reference images
$\{I_{j \rightarrow 0} \}_{j=1}^{n-1}$ with a binary validity mask $\{M_j\}_{j=1}^{n-1}$ indicating valid  projected pixels during warping. 
For ideal unoccluded Lambertian surfaces, pixels  in $\{ I_0, \{I_{j \rightarrow 0} \}_{j=1}^{n-1}  \}$ should be photometrically consistent. 
The multi-view photometric loss enforces this:
\begin{equation}
  \vspace{-0.2cm}
\label{eq:mpc_rgb}
% \mathcal{L}_{\text{mpc}}^{\text{RGB}} 
L = \frac{1}{n-1} \sum_{j=1}^{n-1}  \frac  {\left\| (I_{j \rightarrow 0} - I_{\text{0}}) \odot M_j  \right\|_1} {\left\|  M_j \right\|_1}.
\end{equation}




  % However, this formulation assumes constant lighting across views, which 
% is often violated in practice.
% ===== Illumination-robust Feature Matching =====
\noindent\textbf{Illumination-robust Feature Matching.}
Equation~\ref{eq:mpc_rgb} relies on RGB consistency, which is fragile to 
lighting variations, shadows, and specular reflections common in real scenes. 
To achieve robust geometric supervision, we replace it with feature-based matching using a frozen pre-trained feature network in~\cite{gu2020cascade}:
\begin{equation}
\label{eq:mpc_feat}
\vspace{-0.2cm}
% \mathcal{L}_{\text{mpc}}^{\text{Fea}} 
L = \frac{1}{n-1} \sum_{j=1}^{n-1}  \frac  {\left\| \frac{1}{2}\left(1 - \cos(\mathcal{F}_{\text{0}}, \mathcal{F}_{j \rightarrow 0 } ) \right) \odot M_j  \right\|_1} {\left\|  M_j \right\|_1},
\end{equation}
where $\mathcal{F}_{\text{0}}$ and ${ \mathcal{F}_{j \rightarrow 0 } }$ are features extracted from the reference view and features warped from source views $j$ to the reference view $0$. 
Since features are computed once during preprocessing and remain frozen during training, this incurs negligible computational overhead while significantly improving robustness to illumination changes.


% ===== Occlusion-aware Regularization =====
\noindent\textbf{Occlusion-aware Photometric Consistency.}
Occlusions cause photometric consistency to \textbf{fail} even with accurate depth. 
We address this via \emph{pixel-wise top-k selection}, which adaptively chooses the most reliable correspondences from visible source views.
For each reference pixel $p$, we identify the top-$k$ most consistent correspondences across all warped source features $\{  \mathcal{F}_{j \rightarrow 0 }  \}_{j=1}^{n-1}$:
% \begin{equation}
% \mathcal{T}_k(p) = \underset{S \subset \{1,\ldots,n-1\}, |S|=k}{\arg\min} 
%     \sum_{j \in S} \left\|    \frac{1}{2}\left(1 - \cos(\mathcal{F}_{\text{0}}, \mathcal{F}_{j \rightarrow 0 } )   \right\|_1.
% \end{equation}
\begin{equation}
  \begin{aligned}
  \mathcal{T}_k(p) = \underset{\substack{S \subset \{1,\ldots,n-1\} \\ |S|=k}}{\arg\min} 
      \sum_{j \in S} \left\| \frac{1}{2}\left(1 - \cos(\mathcal{F}_{\text{0}}(p), \mathcal{F}_{j \rightarrow 0}(p))  \right) \right\|_1.
  \end{aligned}
  \end{equation}
This reformulates \cref{eq:mpc_feat} as an adaptive aggregation over $\mathcal{T}_k(p)$:
\begin{equation}
\label{eq:mpc_final}
\mathcal{L}_{\text{mpc}}^{\text{Fea}}(p) = \frac{1}{k} \sum_{j \in \mathcal{T}_k(p)} \left\| \frac{1}{2} \left(1 - \cos(\mathcal{F}_{\text{0}}(p), \mathcal{F}_{j \rightarrow 0}(p)) \right) \right\|_1.
\end{equation}
This pixel-wise selection naturally handles spatially-varying occlusions: for pixels occluded in half the views, the remaining visible views still provide valid supervision. 
Setting $k = \lceil (n-1)/2 \rceil$ balances coverage and outlier rejection.


\subsubsection{Edge-aware Depth Smoothness.}
\label{subsubsec:smooth}

% ===== Edge-aware Smoothness =====
% \noindent\textbf{Edge-aware Depth Smoothness.}
In regions visible from only one views, multi-view photometric consistency fails to provide sufficient geometric constraints. 
We therefore regularize these under-constrained areas with edge-aware depth smoothness:
\begin{equation}
  \vspace{-0.2cm}
\label{eq:smooth}
\mathcal{L}_{\text{smooth}} = \sum_{p} \left\| \nabla D_{\text{0}}(p) \right\|_1 \cdot \exp\left( -\alpha \left\| \nabla I_{\text{0}}(p) \right\|_1 \right),
\end{equation}
where $\nabla D$ and $\nabla I$ denote the depth and image gradients, respectively, and $\alpha=1$ controls edge sensitivity. This encourages smooth depth in textureless regions while preserving discontinuities at object boundaries.

By enforcing geometric consistency through this robust regularization, our 
method significantly improves novel view synthesis quality in weakly-textured regions—
areas where 3D Gaussian Splatting typically struggles due to insufficient photometric 
constraints.
The geometric regularization provides reliable supervision even when 
appearance cues are ambiguous.










\subsection{Geometry-guided Appearance Optimization }
\label{subsec:appearance}
We propose to leverage regularized geometry for appearance optimization through 
virtual-view sampling. 
Existing methods~\cite{zhu2024fsgs,li2024dngaussian,xu2024mvpgs,han2024binocular} 
render virtual novel views to mitigate texture under-constraint in sparse settings, 
enforcing regularization via monocular or binocular depth consistency.
Yet they 
are limited by: scale ambiguity in monocular depth~\cite{zhu2024fsgs,li2024dngaussian}, 
noise in MVS priors~\cite{xu2024mvpgs}, and restricted diversity with inaccurate 
rendered depth~\cite{han2024binocular}.

We instead utilize our regularized Gaussian geometry to enable flexible, reliable 
virtual-view supervision. Our approach comprises: cycle-consistency 
filtering to identify valid depth regions (\cref{sec:depth_filter}), and 
appearance supervision via virtual-view photometric consistency over these 
validated regions (\cref{sec:reliable_synthesis}).

% However, they suffer from key limitations: monocular depth contains scale ambiguity~\cite{zhu2024fsgs,li2024dngaussian}, imprefect pre-trained MVS introduces inevitable noise~\cite{xu2024mvpgs}, and stereo pairs provide limited diversity and suffers inaccurate depth~\cite{han2024binocular}.
% In contrast, we leverage the regularized depth maps rendered from the Gaussian geometry to construct valid supervision with more flexible virtual view sampling.
% This consists of two stages: (\romannumeral1) cycle consistency depth filtering that identifies trustworthy regions (\cref{sec:depth_filter}), 
% and (\romannumeral2) virtual-view appearance supervision via virtual-view photometric consistency from valid regions (\cref{sec:reliable_synthesis}).


% Naively warping training 
% images using rendered depth—as in prior work~\cite{han2024binocular}—suffers 
% from a circular dependency: unreliable depth produces misaligned virtual views, 
% which provide corrupted supervision that further degrades geometry. 
% We address 
% this by restricting it to \emph{geometrically reliable} regions 
% validated through cycle consistency filtering. 


% This ensures that 
% the synthesized virtual parts are coherent with true scene appearance, providing clean 
% supervision without introducing artifacts.

% Our pipeline consists of two stages: (\romannumeral1) cycle consistency depth filtering that identifies trustworthy regions (\cref{sec:depth_filter}), 
% and (\romannumeral2) flexible view synthesis that generates virtual views 
% from valid regions (\cref{sec:reliable_synthesis}).








% ===== 3.3.1 Depth Reliability Filtering =====
\subsubsection{Cycle Consistency Depth Filtering}
\label{sec:depth_filter}
To ensure geometry reliability, we validate rendered depth through cycle-consistency filtering before synthesizing virtual views.
Given a reference view $I_{\text{0}}$, source views $\{I_j\}_{j=1}^{n-1}$, camera 
intrinsic $K$, relative transformations $\{T_{0 \rightarrow j} \}_{j=1}^{n-1}$, and rendered depth 
maps $\{D_i\}_{i=0}^{n-1}$ from Gaussian splatting, we perform forward-backward 
warping.
For each pixel $p$ in $I_{\text{0}}$, we first forward warp to source 
view $I_j$ using depth $D_0(p)$ to obtain projected pixel $p'_{j}$ (see \cref{eq:inverse_warp}). 
We then backward warp $p'_{j}$ to $I_{\text{0}}$ using source depth $D_j(p'_{j})$ 
to obtain reprojected pixel $p''_{j}$ and depth $\tilde{D}_j(p)$:
\begin{equation}
  \vspace{-0.2cm}
\label{eq:inverse_warp_back}
p''_j = KT^{-1}_{0 \rightarrow j}D_j(p'_{j})K^{-1}p'_{j}.
\end{equation}
The depth error between original and reprojected depth measures geometric consistency:
\begin{equation}
\label{eq:depth_error}
\vspace{-0.2cm}
e_j(p) = \left| D_{\text{0}}(p) - \tilde{D}_j(p) \right|.
\end{equation}
A pixel is reliable if its depth error falls below threshold $\tau_d$ for at 
least $m$ of the $n-1$ source views:
\begin{equation}
  \vspace{-0.2cm}
\label{eq:depth_mask}
\mathcal{M}_{\text{reliable}}(p) = \mathbb{I}\left[ \sum_{j=1}^{n-1} \mathbb{I}[e_j(p) < \tau_d] \geq m \right],
\end{equation}
where $m = \lceil (n-1)/2 \rceil$ ensures consistency with at least half the 
sources, and $\tau_d = 0.01 \cdot \max(D_{\text{0}})$. This binary mask 
$\mathcal{M}_{\text{reliable}}$ identifies regions where rendered depth $D_0$ is 
validated by cycle consistency, ensuring subsequent warping produces views 
aligned with true scene structure.


% ===== 3.3.2 Reliable View Synthesis =====
\subsubsection{Virtual-view Photometric Consistency}
\label{sec:reliable_synthesis}
With reliable depth identified, we propagate accurate geometry to unseen appearances via virtual-view photometric consistency. 
Unlike prior stereo-pair approaches~\cite{han2024binocular}, we sample virtual poses $\{\mathcal{P}_v\}_{v=1}^{N_v}$ across a wider range: for each reference position $\mathbf{x}$, we randomly sample within a sphere of radius $r$, providing sufficient viewpoint diversity.

For each virtual view, we forward warp (\cref{eq:inverse_warp}) pixels from all training images $\{I_i\}_{i=0}^{n-1}$ using masked depths $\{\mathcal{M}_{\text{reliable}}^{i} \odot D_i\}_{i=0}^{n-1}$ to synthesize virtual image $\mathcal{I}_{v}$ with validity mask $M_v$, excluding unreliable regions to prevent geometric errors from contaminating supervision.



\begin{table}[t]
  \centering
  \caption{\textbf{Quantitative comparisons on LLFF~\cite{mildenhall2019local} dataset under sparse view settings. }}
  \label{tab:llff_results}
  \resizebox{1.0\linewidth}{!}{
  \begin{tabular}{l|ccc|ccc|ccc}
  \toprule
  \multirow{2}{*}{Methods} & \multicolumn{3}{c}{PSNR↑} & \multicolumn{3}{c}{SSIM↑} & \multicolumn{3}{c}{LPIPS↓} \\

  \cmidrule{2-10}
  & 3-view & 6-view & 9-view & 3-view & 6-view & 9-view & 3-view & 6-view & 9-view \\
  \midrule
  DietNeRF~\cite{jain2021dietnerf} & 14.94 & 21.75 & 24.28 & 0.370 & 0.717 & 0.801 & 0.496 & 0.248 & 0.183 \\
  RegNeRF~\cite{niemeyer2022regnerf} & 19.08 & 23.10 & 24.86 & 0.587 & 0.760 & 0.820 & 0.336 & 0.206 & 0.161\\
  FreeNeRF~\cite{yang2023freenerf} & 19.63 & 23.73 & 25.13 & 0.612 & 0.779 & 0.827 & 0.308 & 0.195 & 0.160 \\
  SparseNeRF~\cite{wang2023sparsenerf} & 19.86 & 23.26 & 24.27 & 0.714 & 0.741 & 0.781 & 0.243 & 0.235 & 0.228 \\
  % ReconFusion~\cite{wu2024reconfusion} & 21.34 & 24.25 & 25.21 & 0.724 & 0.815 & 0.848 & 0.203 & 0.152 & 0.134 \\
  % MuRF~\cite{xu2024murf} & 21.26 & 23.54 & 24.66 & 0.722 & 0.796 & 0.836 & 0.245 & 0.199 & 0.164 \\
  \midrule
  3DGS~\cite{kerbl20233d} & 15.52 & 19.45 & 21.13 & 0.405 & 0.627 & 0.715 & 0.408 & 0.268 & 0.214 \\
  FSGS~\cite{zhu2024fsgs} & 20.31 & 24.20 & 25.32 & 0.652 & 0.811 & 0.856 & 0.288 & 0.173 & 0.136 \\
  DNGaussian~\cite{li2024dngaussian} & 19.12 & 22.18 & 23.17 & 0.591 & 0.755 & 0.788 & 0.294 & 0.198 & 0.180\\
  CoR-GS~\cite{zhang2024cor} & 20.45 & 24.49 & 26.06 & 0.712 & 0.837 & 0.874 & 0.196 & 0.115 & \colorbox{thirdplace}{0.089} \\
  BinocularGS~\cite{han2024binocular} & \colorbox{secondplace}{21.44} & \colorbox{thirdplace}{24.87} & 26.17 & \colorbox{secondplace}{0.751} & \colorbox{thirdplace}{0.845} & \colorbox{thirdplace}{0.877} & \colorbox{secondplace}{0.168} & \colorbox{firstplace}{0.106} & 0.090 \\
  DropGaussians~\cite{park2025dropgaussian} & 20.76 & 24.74 & \colorbox{thirdplace}{26.21} & 0.713 & 0.837 & 0.874 & 0.200 & 0.117 & \colorbox{secondplace}{0.088} \\
  NexusGS~\cite{zheng2025nexusgs} & 21.07 & - & - & 0.738 & - & - & \colorbox{thirdplace}{0.177} & - & - \\
  ComapGS~\cite{jang2025comapgs} & \colorbox{thirdplace}{21.11} & \colorbox{secondplace}{25.20} & \colorbox{firstplace}{26.73} & \colorbox{thirdplace}{0.747} & \colorbox{secondplace}{0.854} & \colorbox{firstplace}{0.886} & 0.182 & \colorbox{secondplace}{0.108} & \colorbox{firstplace}{0.082} \\
  \midrule
  \textbf{Ours} & \colorbox{firstplace}{\textbf{22.20}} & \colorbox{firstplace}{\textbf{25.37}} & \colorbox{secondplace}{\textbf{26.45}} & \colorbox{firstplace}{\textbf{0.778}} & \colorbox{firstplace}{\textbf{0.856}} & \colorbox{secondplace}{\textbf{0.881}} & \colorbox{firstplace}{\textbf{0.157}} & \colorbox{thirdplace}{\textbf{0.109}} & {\textbf{0.096}} \\
  \bottomrule
  \end{tabular}
  }
  % \vspace{-0.2cm}
  \end{table}


\begin{figure*}[t]
\vspace{-0.2cm}
\begin{center}
   \includegraphics[trim={2cm 0cm 2cm 0cm},clip,width=0.72\linewidth]{figures/LLFF_com.pdf}
\end{center}
\vspace{-1cm}
% \caption{ \textbf{Motivation.} }
\caption{\textbf{Visual comparison on LLFF~\cite{mildenhall2019local} dataset.}}
% \label{fig:motivation}
\label{fig:llffcom}
\vspace{-0.3cm}
\end{figure*}








% ===== Augmented Photometric Loss =====
\noindent\textbf{Vitural-view Photometric Consistency Loss.}
The synthesized virtual images are incorporated into training for appearance optimization. 
For each virtual view, we render the Gaussian to acquire the rendered virtual image $\mathcal{I}_{v}^{R}$ from the 
virtual pose $\mathcal{P}_v$ and enforce photometric consistency on valid pixels:
\begin{equation}
  \vspace{-0.2cm}
\label{eq:aug_loss}
L_{app} = \sum_{p \in \mathcal{M}_{\text{v}}}   \left\| \mathcal{I}_{v}(p) - \mathcal{I}_{v}^{R}(p) \right\|_1.
\end{equation}
This serves two purposes: it provides additional observations to optimize appearance in unseen views and prevent overfitting, and conversely, it constrains geometry through supervision from novel viewpoints. 
Critically, because virtual-view images are synthesized from reliability-filtered depth, they provide clean supervision without the geometric distortions introduced by prior methods relying on unreliable depth predictions.









\subsection{Overall Pipeline}
\label{subsec:optimization}
We integrate geometric regularization and geometry-guided appearance 
optimization via curriculum learning~\cite{bengio2009curriculum}.

\noindent\textbf{Training Objective.} The complete loss combines four terms:
  \begin{equation}
    % \vspace{-0.2cm}
    \begin{split}
    \mathcal{L}_{\text{total}} = 
    \mathcal{L}_{\text{3DGS}} 
    &+ \mathcal{L}_{\text{consis}} 
    + \lambda_{\text{mpc}} \mathcal{L}_{\text{mpc}}^{\text{Fea}} \\
    &+ \lambda_{\text{smooth}} \mathcal{L}_{\text{smooth}} 
    + \lambda_{\text{app}} \mathcal{L}_{\text{app}},
    \end{split}
    \end{equation} 
    where $\mathcal{L}_{\text{3DGS}}$ is the base photometric loss, 
    $\mathcal{L}_{\text{consis}}$ enforces binocular consistency inherited from the baseline~\cite{han2024binocular}, 
    $\mathcal{L}_{\text{mpc}}^{\text{Fea}} + \lambda_{\text{smooth}}\mathcal{L}_{\text{smooth}}$ enforces geometry regularization (\cref{subsec:geometry}), 
    and $\mathcal{L}_{\text{app}}$ optimizes appearance in unseen views (\cref{subsec:appearance}).

\noindent\textbf{Optimization.}
We employ a three-stage curriculum: Stage 1 optimizes $\mathcal{L}_{\text{3DGS}}$ to establish coarse geometry; Stage 2  activates geometric regularization $\lambda_{\text{mpc}}  \mathcal{L}_{\text{mpc}}^{\text{Fea}} + \lambda_{\text{smooth}}\mathcal{L}_{\text{smooth}}$; 
Stage 3 adds appearance supervision $\lambda_{\text{app}} \mathcal{L}_{\text{app}}$ from virtual views. 
This staged approach ensures stable convergence.





